% ============================================================
%  Introduction
% ============================================================
\section{Introduction}
\label{sec:intro}

Vision--language models~\cite{radford2021clip,li2022blip} are now the standard
instrument for assembling image--text corpora: a candidate pool is encoded, every
pair receives a similarity, and pairs below a threshold are discarded or queued for
relabelling. This presumes that a similarity can be compared \emph{across} pairs,
that a pair scoring $0.33$ is noisier than one scoring $0.25$. The premise fails.
Figure~\ref{fig:motivation} shows a correctly matched pair scoring $0.252$ and a
pair whose caption misstates the number of animals scoring $0.331$ under the same
frozen encoder.

Existing detectors intervene on the encoder or on the neighbourhood. CSLS
re-normalises a score against a neighbour set, matching heads replace it with a
learned score, and LEMoN judges a pair by its agreement with multimodal neighbours,
while TRACED edits the caption and reads the resulting similarity trajectory. Each improves the signal that is subsequently
thresholded, and each leaves one question unanswered: for a fixed signal, do two
unrelated pairs share a scale? Capacity does not settle it within the zero-shot CLIP
family: ViT-L/14, with $427.6$M parameters against $149.6$M, leaves Attribute AUROC
at $0.6779$ against $0.6806$. The deficiency is not the representation but the
absence of an explicit model of the score component that should not enter a global
comparison.

We therefore model the additively predictable part of the standardised similarity
as $g_\theta=a_\phi(u)+b_\psi(v)+b_0$, fit it to each pair's own score without using
a label, and subtract it at inference, $z'=z-\lambda g_\theta$. The
encoder is frozen; inference consumes one pair, with no reference bank, partner
caption or cycle; labels enter only the selection of $\lambda$. We evaluate on
FGBench, $88{,}024$ reciprocal cycles over five factual facets derived from COCO and
Visual Genome.

\needspace{5\baselineskip}
\noindent\textbf{Contributions.}
\begin{itemize}
  \item \textbf{Method.} CRD, a post-hoc calibration that estimates additive
  image- and text-conditioned score components with two lightweight heads over a
  frozen encoder and subtracts them from the standardised similarity, improving
  global discrimination without any reference structure at inference.
  \item \textbf{Benchmark.} FGBench, an $88{,}024$-cycle reciprocal benchmark with
  five factual facets, explicit split hygiene and an automatic split/leakage audit.
  \item \textbf{Analysis.} Paired tests against five same-backbone detectors and two
  parameter-free debiasing baselines, ablations, selected frozen backbones, and an
  explicit account of where the calibration stops helping.
\end{itemize}
